\documentclass[12pt, a4paper]{article}

% --- Pacotes Fundamentais ---
\usepackage[utf8]{inputenc}
\usepackage[T1]{fontenc}
\usepackage[brazil]{babel}
\usepackage{tikz}
\usetikzlibrary{shapes.geometric, arrows.meta, positioning, calc, fit, backgrounds}

\usepackage[
    a4paper,
    left=3cm,
    right=2cm,
    top=3cm,
    bottom=2cm
]{geometry}

\usepackage{graphicx}
\usepackage{float}
\usepackage{hyperref}
\usepackage{caption}
\usepackage{setspace}
\usepackage{indentfirst} % Indenta o primeiro parágrafo de cada seção
\usepackage{titlesec}    % Formatação de títulos
\usepackage{booktabs}    % Tabelas
\usepackage{tabularx}    % Tabelas com largura ajustada
\usepackage{listings}    % Blocos de código-fonte
\usepackage{xcolor}      % Cores para código e detalhes visuais
\usepackage{colortbl}    % Preenchimento e cores em tabelas
\usepackage{fancyhdr}    % Numeração de páginas no cabeçalho superior direito
\usepackage{amsmath}
\usepackage{enumitem}

% --- Configuração de Cores para Código ---
\definecolor{codegreen}{rgb}{0.15,0.55,0.15}
\definecolor{codegray}{rgb}{0.45,0.45,0.45}
\definecolor{codepurple}{rgb}{0.50,0.10,0.60}
\definecolor{backcolour}{rgb}{0.97,0.97,0.98}
\definecolor{linegray}{rgb}{0.85,0.85,0.85}

\lstdefinestyle{estilocodigo}{
    backgroundcolor=\color{backcolour},   
    commentstyle=\color{codegreen}\itshape,
    keywordstyle=\color{blue}\bfseries,
    numberstyle=\tiny\color{codegray},
    stringstyle=\color{codepurple},
    basicstyle=\ttfamily\footnotesize,
    breakatwhitespace=false,         
    breaklines=true,                 
    captionpos=t,                    
    keepspaces=true,                 
    numbers=left,                    
    numbersep=7pt,                  
    showspaces=false,                
    showstringspaces=false,
    showtabs=false,                  
    tabsize=2,
    frame=single,
    rulecolor=\color{linegray}
}
\lstset{style=estilocodigo}

% --- Formatação Acadêmica ABNT ---
\onehalfspacing
\setlength{\parindent}{1.25cm}
\setlength{\parskip}{0pt}

% Espaçamento de figuras e tabelas
\setlength{\textfloatsep}{12pt plus 2pt minus 2pt}
\setlength{\floatsep}{8pt plus 2pt minus 2pt}
\setlength{\intextsep}{10pt plus 2pt minus 2pt}

% Espaçamento das legendas
\setlength{\abovecaptionskip}{4pt}
\setlength{\belowcaptionskip}{4pt}
\captionsetup{font=small,justification=centering}

% --- Formatação de Títulos ---
\titleformat{\section}
  {\normalfont\normalsize\bfseries}
  {\thesection}{1em}{\MakeUppercase}

\titleformat{\subsection}
  {\normalfont\normalsize\bfseries}
  {\thesubsection}{1em}{}

\titleformat{\subsubsection}
  {\normalfont\normalsize\itshape}
  {\thesubsubsection}{1em}{}

% Configuração de Links
\hypersetup{
    colorlinks=true,
    linkcolor=black,
    citecolor=black,
    urlcolor=blue!70!black
}

% Cabeçalhos e Rodapés
\pagestyle{fancy}
\fancyhf{}
\renewcommand{\headrulewidth}{0pt}
\rhead{\small\thepage}

\begin{document}

% =========================================================================
% ELEMENTOS PRÉ-TEXTUAIS
% =========================================================================

% --- Capa (IFNMG) ---
\begin{titlepage}
    \begin{center}
        \vspace*{-1.2cm}
        \noindent\hspace*{-2cm}%
        \begin{minipage}[c]{3.2cm}
            \centering
            \begin{tikzpicture}[scale=0.48]
                % Linha 4 (topo): circulo + 2 quadrados
                \draw[fill=red!80!black, draw=none] (0,3) circle (0.45);
                \draw[fill=green!60!black, draw=none, rounded corners=1.2pt] (0.55,2.55) rectangle (1.45,3.45);
                \draw[fill=green!60!black, draw=none, rounded corners=1.2pt] (1.55,2.55) rectangle (2.45,3.45);

                % Linha 3: 2 quadrados
                \draw[fill=green!60!black, draw=none, rounded corners=1.2pt] (-0.45,1.55) rectangle (0.45,2.45);
                \draw[fill=green!60!black, draw=none, rounded corners=1.2pt] (0.55,1.55) rectangle (1.45,2.45);

                % Linha 2: 3 quadrados
                \draw[fill=green!60!black, draw=none, rounded corners=1.2pt] (-0.45,0.55) rectangle (0.45,1.45);
                \draw[fill=green!60!black, draw=none, rounded corners=1.2pt] (0.55,0.55) rectangle (1.45,1.45);
                \draw[fill=green!60!black, draw=none, rounded corners=1.2pt] (1.55,0.55) rectangle (2.45,1.45);

                % Linha 1: 2 quadrados
                \draw[fill=green!60!black, draw=none, rounded corners=1.2pt] (-0.45,-0.45) rectangle (0.45,0.45);
                \draw[fill=green!60!black, draw=none, rounded corners=1.2pt] (0.55,-0.45) rectangle (1.45,0.45);
            \end{tikzpicture}
        \end{minipage}%
        \begin{minipage}[c]{\dimexpr\textwidth-1.2cm\relax}
            \centering
            \fontsize{9}{11}\selectfont
            MINISTÉRIO DA EDUCAÇÃO \\
            SECRETARIA DE EDUCAÇÃO PROFISSIONAL E TECNOLÓGICA \\
            INSTITUTO FEDERAL DO NORTE DE MINAS GERAIS -- IFNMG \\
            \small Criado pela Lei nº 11.892, de 29/12/2008 \\
            Campus Montes Claros
        \end{minipage}

        \vfill

        \normalsize
        ANDRÉ FELIPE OLIVEIRA LOPES \\
        JOÃO KENNEDY SOUZA SOARES \\

        \vfill

        \textbf{\large RELATÓRIO TÉCNICO: SISTEMAS DISTRIBUÍDOS}\\[0.35cm]
        \textbf{\normalsize ESTUDO COMPARATIVO ENTRE SOCKETS TCP, RPC (gRPC) E JAVA RMI} \\

        \vfill

        \hfill
        \begin{minipage}{0.52\textwidth}
            \setstretch{1.0}
            \small
            Relatório técnico apresentado como parte das exigências de avaliação da disciplina de Sistemas Distribuídos do Curso de Bacharelado em Ciência da Computação do Instituto Federal do Norte de Minas Gerais -- IFNMG, Campus Montes Claros. \\
            
            \vspace{0.3cm}
            \textbf{Atividade:} Atividade 3 -- Sockets, RPC e RMI \\
            \textbf{Docente:} Prof. Me. Danilo Teixeira Silva
        \end{minipage}

        \vfill

        \normalsize
        MONTES CLAROS - MG \\
        2026
    \end{center}
\end{titlepage}

% Início da contagemdas páginas
\pagestyle{fancy}

% =========================================================================
% 1. INTRODUÇÃO
% =========================================================================
\section{Introdução}
Este trabalho apresenta a implementação e a análise comparativa de três paradigmas fundamentais de comunicação em sistemas distribuídos: Sockets TCP, Remote Procedure Call (gRPC) e Remote Method Invocation (Java RMI). Com o intuito de confrontar empiricamente aspectos como nível de abstração, complexidade de implementação, segurança de tipos e sobrecarga de protocolo, foi desenvolvida uma mesma aplicação para os três cenários: uma calculadora distribuída interativa capaz de realizar as quatro operações aritméticas básicas (soma, subtração, multiplicação e divisão com tratamento de divisão por zero). A padronização da regra de negócio permitiu isolar as características de cada middleware, evidenciando as diferenças reais entre o controle em baixo nível e a transparência de chamadas remotas.

% =========================================================================
% 2. DESCRIÇÃO DA APLICAÇÃO: CALCULADORA DISTRIBUÍDA
% =========================================================================
\section{Descrição da Aplicação: Calculadora Distribuída}

Como base de comparação uniforme entre os três paradigmas, implementou-se uma \textbf{Calculadora Distribuída} sob o modelo cliente-servidor interativo. A aplicação disponibiliza as quatro operações aritméticas básicas sobre valores \texttt{double}: soma ($a + b$), subtração ($a - b$), multiplicação ($a \times b$) e divisão ($a \div b$), incluindo validação estrita contra divisão por zero.

O cliente opera via console em laço contínuo através da sintaxe \texttt{<OPERACAO> <NUM1> <NUM2>} (ex.: \texttt{SOMA 10 20}), permanecendo ativo até o comando \texttt{SAIR}. O servidor computa a operação e retorna o valor numérico calculado ou uma mensagem de erro correspondente. A Figura~\ref{fig:fluxo_geral} ilustra o fluxo de comunicação.

\begin{figure}[H]
    \centering
    \begin{tikzpicture}[
        node distance=5.6cm,
        box/.style={
            rectangle, 
            draw=black, 
            thick, 
            fill=blue!5, 
            rounded corners=3pt, 
            minimum height=1.6cm, 
            minimum width=3.8cm, 
            align=center
        },
        arrow/.style={-{Stealth[length=2.5mm]}, thick}
    ]
        \node[box] (cli) {\textbf{Cliente}\\(Console no terminal)};
        \node[box, right=of cli] (srv) {\textbf{Servidor}\\(Execução)};

        \draw[arrow] ([yshift=0.3cm]cli.east) -- node[midway, above=2pt, font=\footnotesize] {Operação e Parâmetros} ([yshift=0.3cm]srv.west);
        \draw[arrow] ([yshift=-0.3cm]srv.west) -- node[midway, below=2pt, font=\footnotesize] {Valor Calculado ou Erro} ([yshift=-0.3cm]cli.east);
    \end{tikzpicture}
    \caption{Fluxo de comunicação cliente-servidor da Calculadora Distribuída.}
    \label{fig:fluxo_geral}
\end{figure}

% =========================================================================
% 3. ARQUITETURA E IMPLEMENTAÇÃO DAS ABORDAGENS
% =========================================================================
\section{Arquitetura e Implementação das Abordagens}

A seguir, apresentam-se as características arquiteturais, o funcionamento e os diagramas de sequência de cada tecnologia.

% -------------------------------------------------------------------------
\subsection{Sockets TCP}

A implementação opera na camada de transporte via fluxos confiáveis orientados a conexão (\texttt{ServerSocket} e \texttt{Socket}) na porta \texttt{5000}. O servidor adota o padrão \textit{thread-per-connection} através da classe \texttt{AtendenteCliente}, atendendo clientes concorrentes de forma isolada sem bloquear o laço de escuta.

Por se tratar de um canal de fluxo de bytes não estruturado, definiu-se um protocolo textual delimitado por quebra de linha (\texttt{\textbackslash n}). A validação e o \textit{parsing} dos comandos ocorrem manualmente via \texttt{String.split()} e \texttt{Double.parseDouble()}, retornando o resultado ou mensagens de falha (\texttt{ERRO: ...}) sem derrubar a conexão. A Figura~\ref{fig:seq_socket} detalha esse fluxo.

\begin{figure}[H]
    \centering
    \includegraphics[width=0.85\textwidth]{diagrama-sequencia-socket.png}
    \caption{Diagrama de sequência da interação cliente-servidor via Sockets TCP.}
    \label{fig:seq_socket}
\end{figure}

% -------------------------------------------------------------------------
\subsection{RPC com gRPC}

A versão RPC utiliza o arcabouço gRPC sobre HTTP/2 com serialização binária compacta em Protocol Buffers (proto3). O contrato das operações e mensagens foi formalizado no arquivo \texttt{calculadora.proto} (Listagem~\ref{lst:proto}).

\begin{lstlisting}[language=Java, caption={Contrato de serviço definido em Protocol Buffers (calculadora.proto).}, label={lst:proto}]
syntax = "proto3";
option java_multiple_files = true;
option java_package = "org.ifnmg.rpc";

service CalculadoraService {
  rpc Somar (OperacaoRequest) returns (ResultadoResponse);
  rpc Subtrair (OperacaoRequest) returns (ResultadoResponse);
  rpc Multiplicar (OperacaoRequest) returns (ResultadoResponse);
  rpc Dividir (OperacaoRequest) returns (ResultadoResponse);
}

message OperacaoRequest {
  double primeiro = 1;
  double segundo = 2;
}

message ResultadoResponse {
  double resultado = 1;
}
\end{lstlisting}

O compilador \texttt{protobuf-maven-plugin} gera as classes de dados tipadas e o serviço base (\texttt{CalculadoraServiceImplBase}). O servidor (\texttt{ServidorRpc}) executa sobre o Netty na porta \texttt{8080}, enquanto o cliente (\texttt{ClienteRpc}) utiliza um canal gerenciado (\texttt{ManagedChannel}) com \textit{stub} síncrono bloqueante (\texttt{BlockingStub}), simplificando a invocação como ilustrado na Figura~\ref{fig:seq_grpc}.

\begin{figure}[H]
    \centering
    \includegraphics[width=1\textwidth]{diagrama-sequencia-rpc.png}
    \caption{Diagrama de sequência da invocação remota com gRPC e Protocol Buffers.}
    \label{fig:seq_grpc}
\end{figure}

% -------------------------------------------------------------------------
\subsection{Java RMI}

O Java RMI provê computação distribuída baseada em objetos nativos da JVM sobre o protocolo JRMP. A especificação utiliza a interface remota \texttt{CalculadoraRmi} (que estende \texttt{Remote} e declara \texttt{RemoteException}) e a classe concreta \texttt{CalculadoraRmiImpl} (que herda de \texttt{UnicastRemoteObject}).

O nó servidor publica a instância remota no \textbf{RMI Registry} (porta 1099) sob o identificador \texttt{"CalculadoraService"} através do método \texttt{rebind()}. O cliente obtém a referência remota via \texttt{lookup()} e invoca os métodos diretamente sobre o \textit{stub} local com transparência sintática, conforme a Figura~\ref{fig:seq_rmi}.

\begin{figure}[H]
    \centering
    \includegraphics[width=1\textwidth]{diagrama-sequencia-rmi.png}
    \caption{Diagrama de sequência da resolução e chamada remota via Java RMI.}
    \label{fig:seq_rmi}
\end{figure}

% =========================================================================
% 4. ANÁLISE COMPARATIVA DETALHADA
% =========================================================================
\section{Análise Comparativa}

Nesta seção, estabelece-se um comparativo técnico e analítico entre as três soluções com respeito a abstração, arquitetura de rede, complexidade e desempenho.

% -------------------------------------------------------------------------
\subsection{Nível de Abstração Oferecido}

As três soluções situam-se em patamares substancialmente distintos ao longo da escala de abstração, conforme sintetizado na Figura~\ref{fig:abstracao}:

\begin{figure}[H]
    \centering
    \begin{tikzpicture}[
        scale=0.92, every node/.style={transform shape},
        levelbox/.style={rectangle, draw=black!80, thick, rounded corners=3pt, minimum width=4.3cm, minimum height=2.3cm, align=center, font=\small}
    ]
        \node[levelbox, fill=red!8] (sock) at (0, 0) {
            \textbf{Sockets TCP}\\[1mm]
            \textit{Baixo Nível}\\[1.5mm]
            \scriptsize $\bullet$ Fluxo contínuo de bytes\\
            \scriptsize $\bullet$ Protocolo textual manual\\
            \scriptsize $\bullet$ Gestão manual de threads
        };

        \node[levelbox, fill=blue!8] (rpc) at (5.0, 0) {
            \textbf{gRPC (RPC)}\\[1mm]
            \textit{Médio-Alto Nível}\\[1.5mm]
            \scriptsize $\bullet$ Chamada de procedimentos\\
            \scriptsize $\bullet$ Contrato formal em IDL\\
            \scriptsize $\bullet$ Serialização binária
        };

        \node[levelbox, fill=green!8] (rmi) at (10.0, 0) {
            \textbf{Java RMI}\\[1mm]
            \textit{Alto Nível}\\[1.5mm]
            \scriptsize $\bullet$ Orientação a objetos pura\\
            \scriptsize $\bullet$ Transparência sintática\\
            \scriptsize $\bullet$ Tipagem nativa na JVM
        };

        % Seta horizontal de evolução no rodapé
        \draw[-{Stealth[length=3.5mm]}, thick, draw=blue!80!black] 
            (-2.15, -1.6) -- (12.15, -1.6) 
            node[midway, below=3pt, font=\small\bfseries] {Evolução do Nível de Abstração};
    \end{tikzpicture}
    \caption{Escala de abstração entre os paradigmas de comunicação avaliados.}
    \label{fig:abstracao}
\end{figure}

No modelo de \textbf{Sockets}, o desenvolvedor é compelido a lidar com toda a infraestrutura física de rede: buffers de leitura, controle de fim de linha, conversão de tipos de dados de primitivos para strings e controle do ciclo de vida das conexões. 

No \textbf{gRPC}, a complexidade de transporte é totalmente encapsulada, fornecendo a semântica de chamada de procedimento remoto. Contudo, há uma distinção arquitetural sadia entre objetos locais e mensagens de requisição/resposta (\textit{Data Transfer Objects} com \texttt{OperacaoRequest} e \texttt{ResultadoResponse}), prevenindo a ilusão de que a rede possui custo nulo.

Por fim, o \textbf{Java RMI} atinge o topo da abstração sintática: o objeto remoto é manuseado indistinguivelmente de um objeto local, bastando que métodos sejam invocados sobre a interface. No entanto, essa total transparência de localização pode ocultar latências de rede e potenciais falhas parciais que não ocorrem na memória local.

% -------------------------------------------------------------------------
\subsection{Diferenças de Implementação e Protocolos de Rede}

As divergências estruturais entre as tecnologias refletem diretamente no consumo de rede e no acoplamento técnico. Enquanto os Sockets operam sobre TCP puro com mensagens em texto plano delimitadas por \texttt{\textbackslash n} e exigem controle manual de \textit{threads}, o gRPC emprega conexões persistentes multiplexadas sobre HTTP/2 com codificação binária compacta (Protocol Buffers) e \textit{pool} reativo de eventos via Netty, maximizando a eficiência de transmissão. Por sua vez, o Java RMI utiliza o protocolo proprietário JRMP sobre TCP com gestão de \textit{threads} pela própria JVM, mas introduz maior sobrecarga de dados decorrente da serialização de metadados e assinaturas das classes Java.

% -------------------------------------------------------------------------
\subsection{Facilidade de Uso e Complexidade de Manutenção}

\begin{itemize}
    \item \textbf{Sockets:} Inicialização simples para pequenos protótipos, porém apresenta custo de manutenção exponencial. Qualquer acréscimo funcional (como uma nova operação de potenciação) demanda alterações simultâneas no \textit{parser} do servidor, validação sintática do cliente e na documentação informal do protocolo textual.
    \item \textbf{gRPC:} Exige esforço inicial considerável para aprendizado da sintaxe proto3 e parametrização dos \textit{plugins} do Maven. Em contrapartida, a manutenção a longo prazo é extremamente facilitada: adicionar uma nova função resume-se a declarar um método no arquivo \texttt{.proto} e reexecutar a compilação, cabendo ao compilador apontar inconsistências estáticas em tempo de construção.
    \item \textbf{Java RMI:} Altamente intuitivo para engenheiros com familiaridade em Java puro, dispensando bibliotecas externas. Entretanto, sofre de acoplamento severo: alterações na interface remota exigem que clientes e servidores sincronizem seus arquivos \texttt{.class} e garantam a paridade de versões por meio do \texttt{serialVersionUID}.
\end{itemize}

% -------------------------------------------------------------------------
\subsection{Quadro Comparativo Síntese}

A Tabela~\ref{tab:comparativo} consolida os parâmetros técnicos e arquiteturais observados durante a execução do projeto prático.

\begin{table}[H]
    \centering
    \caption{Quadro comparativo entre Sockets TCP, gRPC e Java RMI.}
    \label{tab:comparativo}
    \small
    \renewcommand{\arraystretch}{1.45}
    \begin{tabularx}{\textwidth}{|
        >{\columncolor{gray!8}\raggedright\arraybackslash\hyphenpenalty=10000}p{3.5cm}|
        >{\raggedright\arraybackslash\hyphenpenalty=10000}X|
        >{\raggedright\arraybackslash\hyphenpenalty=10000}X|
        >{\raggedright\arraybackslash\hyphenpenalty=10000}X|
    }
        \hline
        \rowcolor{gray!20}
        \textbf{Critério de Análise} & \textbf{Sockets TCP} & \textbf{gRPC (RPC)} & \textbf{Java RMI} \\ \hline
        \textbf{Camada de Transporte} & TCP Puro & HTTP/2 sobre TCP & JRMP sobre TCP \\ \hline
        \textbf{Formato de Codificação} & Texto delimitado (\texttt{\textbackslash n}) & Binário (Protocol Buffers) & Binário (Serialização Java) \\ \hline
        \textbf{Definição do Contrato} & Informal / Implícito & Formal (IDL proto3) & Formal (Interface Java) \\ \hline
        \textbf{Gestão de Concorrência} & Manual (\texttt{Thread} dedicada) & Automática (\textit{Pool} Netty) & Automática (\textit{Pool} JVM) \\ \hline
        \textbf{Verificação de Tipos} & Em execução (\textit{parsing}) & Em compilação (Forte) & Em compilação (Forte) \\ \hline
        \textbf{Interoperabilidade} & Universal (qualquer SO) & Universal (Poliglota) & Nula (Exclusivo Java/JVM) \\ \hline
        \textbf{Dependências Externas} & Nenhuma (JDK nativo) & Bibliotecas gRPC / Netty & Nenhuma (JDK nativo) \\ \hline
        \textbf{Sobrecarga de Rede} & Baixa (porém textual) & Mínima (ultra-compacto) & Alta (metadados de classes) \\ \hline
        \textbf{Nível de Abstração} & Baixo Nível & Médio-Alto Nível & Alto Nível \\ \hline
    \end{tabularx}
    \begin{center}
        \footnotesize \textbf{Fonte:} Elaborado pelos autores (2026).
    \end{center}
\end{table}

% =========================================================================
% 5. VANTAGENS E DESVANTAGENS PERCEBIDAS
% =========================================================================
\section{Vantagens e Desvantagens Percebidas}

A partir do desenvolvimento e da execução dos experimentos, identificaram-se os pontos fortes e as limitações de cada tecnologia.

\subsection{Sockets TCP}
\begin{itemize}[label=\textcolor{green!60!black}{\textbf{+}}]
    \item \textbf{Controle Total e Leveza:} Não depende de bibliotecas externas. O desenvolvedor tem controle direto de tudo o que envia, podendo criar mensagens simples sob medida (ex.: o texto enxuto \texttt{"SOMA 10 20\textbackslash n"}).
    \item \textbf{Interoperabilidade Universal:} Como a comunicação trafega apenas bytes pela rede, qualquer linguagem consegue se comunicar facilmente (ex.: um cliente simples em Python ou C conecta diretamente no servidor Java).
\end{itemize}
\begin{itemize}[label=\textcolor{red!80!black}{\textbf{--}}]
    \item \textbf{Muito Código Repetitivo:} É preciso programar tudo manualmente, desde gerenciar conexões e \textit{threads} até controlar leituras com \texttt{BufferedReader}, envios com \texttt{PrintWriter} e quebras de linha com \texttt{\textbackslash n}.
    \item \textbf{Facilidade para Erros de Formatação:} Não há checagem automática dos dados. Se o cliente enviar uma mensagem inesperada (ex.: \texttt{"SOMA 10"} faltando um número), o servidor quebra com erros de conversão caso o desenvolvedor não trate cada detalhe manualmente.
\end{itemize}

\subsection{gRPC (RPC)}
\begin{itemize}[label=\textcolor{green!60!black}{\textbf{+}}]
    \item \textbf{Desempenho e Eficiência:} A multiplexação do HTTP/2 aliada à codificação binária do Protocol Buffers proporciona a menor latência e consumo de banda entre as três opções.
    \item \textbf{Ecossistema Poliglota:} Permite que clientes implementados em Python, C++, Go ou Node.js comuniquem-se nativamente com o servidor Java sem nenhuma camada de conversão adicional.
    \item \textbf{Geração de Código Estável:} Minimiza erros humanos na definição de mensagens de transporte.
\end{itemize}
\begin{itemize}[label=\textcolor{red!80!black}{\textbf{--}}]
    \item \textbf{Curva de Entrada e Dependências:} Requer instalação e parametrização de compiladores de protocolo (\texttt{protoc}) e plugins Maven externos.
    \item \textbf{Depuração Opaca:} Diferente do texto cru transmitido em sockets, a inspeção de pacotes em trânsito (por exemplo, via Wireshark) requer chaves de decodificação e arquivos proto para leitura dos dados binários.
\end{itemize}

\subsection{Java RMI}
\begin{itemize}[label=\textcolor{green!60!black}{\textbf{+}}]
    \item \textbf{Integração Orientada a Objetos:} A fluidez de desenvolvimento é máxima para programadores Java, permitindo passar objetos complexos por parâmetro com suporte a polimorfismo.
    \item \textbf{Nativo no JDK:} Não requer inclusão de dependências externas ou plugins de compilação adicionais.
\end{itemize}
\begin{itemize}[label=\textcolor{red!80!black}{\textbf{--}}]
    \item \textbf{Confinamento de Plataforma:} Incapacidade total de interagir com tecnologias fora do ambiente da Máquina Virtual Java.
    \item \textbf{Sobrecarga de Serialização e Segurança:} A serialização padrão do Java é pesada e historicamente suscetível a vulnerabilidades de injeção de classes quando exposta a redes abertas.
\end{itemize}

% =========================================================================
% 6. DIFICULDADES ENCONTRADAS E SOLUÇÕES ADOTADAS
% =========================================================================
\section{Dificuldades Encontradas e Soluções Adotadas}

Durante o ciclo de desenvolvimento das três soluções, os seguintes desafios técnicos foram superados:

\begin{enumerate}
    \item \textbf{Robustez do Parsing Textual nos Sockets:}
    No modelo inicial em Sockets, entradas com múltiplos espaços consecutivos causavam quebras no método \texttt{split()}. A solução consistiu em empregar a expressão regular \texttt{"\textbackslash\textbackslash s+"}.

    \item \textbf{Integração do Plugin Protobuf no Maven (gRPC):}
    A compilação do gRPC depende de binários nativos do \texttt{protoc} adequados à plataforma do sistema operacional hospedeiro. Para evitar incompatibilidades entre sistemas operacionais dos integrantes da dupla, usou-se o \texttt{os-maven-plugin} acoplado ao \texttt{protobuf-maven-plugin} na raiz do módulo RPC, permitindo a detecção dinâmica de arquiteturas Linux x86\_64 e Windows.

    \item \textbf{Padronização Multi-Módulo com Perfis Maven:}
    Para viabilizar uma experiência de compilação integrada e execução uniforme sem necessidade de comandos extensos no terminal, construiu-se uma estrutura de projeto multi-módulo com um arquivo \texttt{pom.xml} agregador na raiz e perfis padronizados (\texttt{-Pservidor} e \texttt{-Pcliente}) em cada subdiretório.
\end{enumerate}

% =========================================================================
% 7. CASOS DE USO PRÁTICOS NA INDÚSTRIA
% =========================================================================
\section{Casos de Uso Práticos na Indústria}

A seleção do mecanismo de comunicação em projetos reais da engenharia de software é guiada por requisitos funcionais e não-funcionais:

\begin{itemize}
    \item \textbf{Indicações para Sockets TCP:}
    Recomendado para sistemas que exigem latência crítica e controle direto do canal de comunicação, tais como servidores de \textbf{jogos eletrônicos multijogador em tempo real}, gateways de sensores para Internet das Coisas (IoT) industrial, e sistemas embarcados.
    
    \item \textbf{Indicações para gRPC (RPC):}
    Consolidou-se como o padrão de fato para a comunicação interna em \textbf{arquiteturas de microsserviços em nuvem de alta escala} (adotado por empresas como Google, Netflix, Uber e Spotify). É a escolha ideal quando sistemas heterogêneos escritos em diferentes linguagens precisam interoperar com alto rendimento, transmissão contínua de dados (\textit{streaming}) e contratos rígidos.

    \item \textbf{Indicações para Java RMI:}
    Apropriado fundamentalmente para a manutenção e evolução de \textbf{sistemas corporativos legados em Java} (como barramentos baseados em Java EE e Enterprise JavaBeans -- EJB), bem como sistemas distribuídos acadêmicos ou corporativos fechados, onde há certeza absoluta de que todos os nós participantes são integralmente baseados na JVM.
\end{itemize}

% =========================================================================
% 8. GUIA DE EXECUÇÃO E REPRODUTIBILIDADE DOS TESTES
% =========================================================================
% \section{Guia de Execução e Testes}

% Para reproduzir os testes e validar o comportamento idêntico das três implementações da calculadora, deve-se proceder conforme as instruções a seguir a partir de um terminal com suporte ao Apache Maven e JDK 17+:

% \subsection{Compilação Geral}
% Na raiz do repositório clonado:
% \begin{lstlisting}[language=bash, numbers=none]
% mvn clean compile
% \end{lstlisting}

% \subsection{Execução das Aplicações}
% Para cada tecnologia, devem-se abrir dois terminais distintos (um para o Servidor e outro para o Cliente):

% \begin{itemize}
%     \item \textbf{Versão 1: Sockets TCP}
%     \begin{itemize}
%         \item Terminal 1 (Servidor): \texttt{mvn exec:java -pl Socket -Pservidor}
%         \item Terminal 2 (Cliente): \texttt{mvn exec:java -pl Socket -Pcliente}
%     \end{itemize}

%     \item \textbf{Versão 2: gRPC (RPC)}
%     \begin{itemize}
%         \item Terminal 1 (Servidor): \texttt{mvn exec:java -pl RPC -Pservidor}
%         \item Terminal 2 (Cliente): \texttt{mvn exec:java -pl RPC -Pcliente}
%     \end{itemize}

%     \item \textbf{Versão 3: Java RMI}
%     \begin{itemize}
%         \item Terminal 1 (Servidor): \texttt{mvn exec:java -pl RMI -Pservidor}
%         \item Terminal 2 (Cliente): \texttt{mvn exec:java -pl RMI -Pcliente}
%     \end{itemize}
% \end{itemize}

% =========================================================================
% 9. CONCLUSÃO
% =========================================================================
\section{Conclusão}

O desenvolvimento prático das três versões da Calculadora Distribuída mostrou claramente a relação de equilíbrio em sistemas distribuídos: quanto maior o controle em baixo nível, maior o trabalho manual de manutenção; e quanto maior a abstração, mais simples e produtivo se torna o código.

A versão com \textbf{Sockets} deixou clara a dificuldade de cuidar de tudo manualmente. Embora ofereça controle total da rede sem depender de bibliotecas extras, exige muito código repetitivo para gerenciar conexões, criar \textit{threads} e tratar mensagens de texto, sendo muito fácil quebrar com entradas mal formatadas.

O \textbf{Java RMI} mostrou como a orientação a objetos facilita a comunicação distribuída, permitindo chamar métodos remotos como se fossem locais. Por outro lado, revelou duas limitações importantes: funciona exclusivamente dentro do ecossistema Java e carrega o peso da serialização nativa da linguagem.

Por fim, o \textbf{gRPC} provou ser a opção mais equilibrada e alinhada ao mercado atual. Ao combinar um contrato bem definido e compatível com várias linguagens (Protocol Buffers) com a alta velocidade do HTTP/2, ele resolve a falta de tipagem dos sockets e supera o isolamento tecnológico do RMI, destacando-se como a melhor escolha para sistemas distribuídos modernos.

% =========================================================================
% ELEMENTOS PÓS-TEXTUAIS: REFERÊNCIAS BIBLIOGRÁFICAS
% =========================================================================
\newpage
\begin{thebibliography}{99}
    \bibitem{coulouris}
    COULOURIS, George; DOLLIMORE, Jean; KINDBERG, Tim; BLAIR, Gordon. \textbf{Sistemas Distribuídos: Conceitos e Projeto}. 5. ed. Porto Alegre: Bookman, 2013.

    \bibitem{tanenbaum}
    TANENBAUM, Andrew S.; STEEN, Maarten van. \textbf{Sistemas Distribuídos: Princípios e Paradigmas}. 2. ed. São Paulo: Pearson Prentice Hall, 2007.

    \bibitem{oracle_rmi}
    ORACLE. \textbf{Java Remote Method Invocation (RMI) Specification}. Java Platform, Standard Edition Documentation. Disponível em: \url{https://docs.oracle.com/javase/8/docs/platform/rmi/spec/rmiTOC.html}. Acesso em: 24 set. 2026.

\end{thebibliography}

\end{document}
